\section*{Contoh}
\addcontentsline{toc}{section}{Contoh}

\begin{contoh}[Menemukan handle]
Pada bentuk kalimat \(\code{id} + \code{id} * \code{id}\), handle
pertama adalah \code{id} paling kiri. Setelah direduksi menjadi
\(F\), \(T\), dan \(E\), handle berikutnya adalah \code{id} kedua.
Pemilihan handle yang salah, misalnya mereduksi seluruh
\(\code{id} + \code{id}\) terlebih dahulu, menghasilkan pohon yang
tidak sesuai precedence.
\end{contoh}

\begin{contoh}[Relasi precedence]
Untuk \(\code{id} + \code{id} * \code{id}\):
\[
\code{id} \cdot> +, \quad
+ <\cdot \code{id}, \quad
+ <\cdot *, \quad
* \cdot> \code{id}.
\]
Relasi ini memaksa perkalian direduksi sebelum penjumlahan.
\end{contoh}

\begin{contoh}[Unary minus]
Operator unary minus memiliki precedence lebih tinggi daripada
binary minus. Pada ekspresi \(-a * b\), unary minus diterapkan pada
\(a\) terlebih dahulu, lalu hasilnya dikalikan \(b\). Grammar harus
membedakan aturan unary dan binary, atau parser generator harus
mendeklarasikan precedence unary secara eksplisit.
\end{contoh}

\begin{contoh}[Kesalahan pada operator precedence]
Pada masukan \(\code{id} + + \code{id}\), setelah shift \(+\) kedua,
tabel tidak memiliki relasi antara \(+\) dan \(+\). Parser
melaporkan kesalahan pada token kedua. Pemulihan dapat dilakukan
dengan melewati token sampai operator atau operand berikutnya.
\end{contoh}

\begin{contoh}[Membaca keluaran kalkulator]
Program \file{shift\_reduce.c} mencetak urutan shift dan reduce.
Untuk masukan \code{2+3*4}, program mencetak shift angka 2, shift
operator +, shift angka 3, shift operator *, shift angka 4, lalu
reduce 3*4, reduce 2+12, dan hasil 14. Urutan ini menunjukkan bahwa
perkalian direduksi lebih dahulu.
\end{contoh}
